\documentclass{book}
\usepackage[utf8]{inputenc}
\usepackage[T1]{fontenc}
\usepackage{geometry}
\usepackage{graphicx}
\usepackage{float}
\usepackage{hyperref}
\usepackage[english,indonesian]{babel}

\title{Optimasi Pemrosesan Template LaTeX Menggunakan Pendekatan Berbasis AI}
\author{Academic Support Agent}
\date{\today}

\begin{document}
\maketitle

\begin{abstract}

\end{abstract}

\chapter{PENDAHULUAN}


\chapter{TINJAUAN PUSTAKA}


% Original content (commented out):
% \documentclass[12pt,a4paper,twoside]{book}
\usepackage[utf8]{inputenc}
\usepackage[spanish,es-noindentfirst,es-tabla]{babel}
\usepackage[x11names]{xcolor}
\usepackage{csquotes}
\usepackage{hyperref}
% La plantilla se creo por Milton Calderón
% Cualquier duda o comentario mvcalderonc@correo.udistrital.edu.co
%requerido para la cabecera
\usepackage{fancyhdr}
\setlength{\headheight}{15pt}
\usepackage{multirow}

\usepackage{xargs}

\newcommandx{\newpar}{\vspace{12pt}\newline}
\usepackage{geometry}
\geometry{
    a4paper,
    left=30mm,
    right=30mm,
    top=30mm,
    bottom=30mm,
}

%Setting library and location of images
\usepackage{graphicx}
\graphicspath{ {images/} }

% package to use numbered list
\usepackage{enumitem}

%%%%%%%%%%%%%%%%%%%%%%%%%%%%%%%%%%%
% SETTING BIBLIOGRAPHY
%%%%%%%%%%%%%%%%%%%%%%%%%%%%%%%%%%%
%\usepackage{biblatex}
%Setting style
\usepackage[
    sortcites,
    backend=biber,
    style=apa,
  ]{biblatex}
%Setting file with bibliography
\addbibresource{sections/references.bib}

%%%%%%%%%%%%%%%%%%%%%%%%%%%%%%%%%%%%%%%%%%%%%
%   Datos de la tesis
%%%%%%%%%%%%%%%%%%%%%%%%%%%%%%%%%%%%%%%%%%%%%
\def\thesistitle{Título de la tesis}
\def\thesisauthorfirsttitle{Ing.}
\def\thesisauthorfirst{Nombre estudiante}
\def\thesissupervisorfirsttitle{Ing.}
\def\thesissupervisorfirst{nombre director}
\def\thesissupervisorfirsttitlelevel{PhD}
\def\thesisdate{2020}

\title{\thesistitle}
\author{\thesisauthorfirst \and \thesissupervisorfirst}
\date{\thesisdate}


\usepackage[onehalfspacing]{setspace}

% Sólo muestra sección y subsección
\setcounter{tocdepth}{2}

\begin{document}
  \frontmatter
  % portada
  \include{sections/portada}
  % contraportada
  \pagenumbering{gobble}
  \include{sections/contraportada}
  \pagenumbering{roman}
  % dedicatoria
  \include{sections/dedicatoria}
  % agradecimientos
  \include{sections/agradecimientos}
  % resumen
  % ..palabras clave
  \include{sections/resumen}
  % tabla de contenido
  % tabla de figuras
  % lista de tablas
  % lista de abreviaturas
  \renewcommand{\contentsname}{CONTENIDO}
  \tableofcontents
  \addcontentsline{toc}{chapter}{CONTENIDO}
  \renewcommand{\listfigurename}{Lista de figuras}
  \listoffigures
  \addcontentsline{toc}{chapter}{Lista de figuras}
  \renewcommand{\listtablename}{Lista de tablas}
  \listoftables
  \addcontentsline{toc}{chapter}{Lista de tablas}
  \mainmatter
  % intro 1
  \include{chapters/introduccion}
  % marco de referencia 2
  \include{chapters/marcoReferencia}
  % metodologia 3
  \include{chapters/metodologia}
  % resultados 4
  \include{chapters/resultados}
  % conclusiones 5
  \include{chapters/conclusiones}
  \backmatter
  % bibliografía
  \include{sections/bibliografia}

\end{document}


\end{document}